\documentclass[aspectratio=169,11pt]{beamer}

% ── Session 2: "Secrets You Cannot Tell" (tacit knowledge / TKBA) ─────────
% Shared brand theme + \gwtitle/\punch/\memeslot live in gwtheme.tex.

\usepackage[utf8]{inputenc}
\usepackage[T1]{fontenc}
\usepackage{lmodern}
\usepackage{../style/greatwall}
\usepackage{booktabs}
\usepackage{array}
\usepackage{colortbl}

\graphicspath{{./}{../assets/}}
\def\gwsessionno{2}
% gwtheme.tex — shared Beamer theme for the Great Wall pilot-course decks.
% \input this from each slides/sessionN.tex *after* loading greatwall.sty and
% *after* defining \gwsessionno (the session number, e.g. \def\gwsessionno{2}).
% Provides the brand look plus three helpers: \gwtitle, \punch, \memeslot.
%
% Depends on: beamer, greatwall.sty (palette + boxes + marks), tikz.

\providecommand{\gwsessionno}{?}

% ── Minimal custom theme (no external font deps; texlive-full only) ────────
\usefonttheme{professionalfonts}
\setbeamertemplate{navigation symbols}{}
\setbeamercolor{normal text}{fg=darktext, bg=white}
\setbeamercolor{structure}{fg=accent}
\setbeamercolor{frametitle}{fg=white, bg=accent}
\setbeamercolor{title}{fg=white}
\setbeamercolor{block title}{fg=white, bg=accentdark}
\setbeamercolor{block body}{fg=darktext, bg=accentlight}
\setbeamercolor{itemize item}{fg=accent}
\setbeamercolor{itemize subitem}{fg=btccol}
\setbeamerfont{frametitle}{series=\bfseries}
\setbeamerfont{title}{series=\bfseries, size=\Huge}
\setbeamertemplate{itemize item}{\textbf{--}}

% Frametitle as a solid green bar
\setbeamertemplate{frametitle}{%
  \nointerlineskip
  \begin{beamercolorbox}[wd=\paperwidth,ht=2.6ex,dp=1.2ex,leftskip=0.4cm,rightskip=0.4cm]{frametitle}%
    \usebeamerfont{frametitle}\insertframetitle%
  \end{beamercolorbox}%
}

% Footline: course | session | frame number
\setbeamercolor{gwfootline}{fg=graytext, bg=white}
\setbeamertemplate{footline}{%
  \leavevmode%
  \hbox{%
  \begin{beamercolorbox}[wd=.5\paperwidth,ht=2.4ex,dp=1ex,leftskip=0.4cm]{gwfootline}%
    \tiny Great Wall --- Pilot Course%
  \end{beamercolorbox}%
  \begin{beamercolorbox}[wd=.5\paperwidth,ht=2.4ex,dp=1ex,rightskip=0.4cm]{gwfootline}%
    \hfill\tiny Session \gwsessionno{} \;\textbullet\; \insertframenumber%
  \end{beamercolorbox}}%
}

% ── Helpers ───────────────────────────────────────────────────────────────
% Big punchy centered statement frame (full-bleed green).
\newcommand{\punch}[1]{%
  \begin{frame}[plain]
  \begin{tikzpicture}[remember picture,overlay]
    \fill[accent] (current page.south west) rectangle (current page.north east);
  \end{tikzpicture}
  \begin{center}\vfill
  {\color{white}\Huge\bfseries #1}\par
  \vfill\end{center}
  \end{frame}}

% Meme slot — shows the image if present, else a labelled placeholder box.
% #1 = filename in ../assets/, #2 = caption/placeholder text.
\newcommand{\memeslot}[2]{%
  \IfFileExists{../assets/#1}{%
    \includegraphics[width=\linewidth,height=0.62\textheight,keepaspectratio]{#1}%
  }{%
    \begin{tcolorbox}[colback=accentlight,colframe=graytext,sharp corners,
      halign=center, valign=center, height=0.55\textheight]
      \color{graytext}\itshape [meme slot: \texttt{#1}]\\[4pt]\small #2
    \end{tcolorbox}%
  }%
}

% Branded title frame. #1 = title, #2 = subtitle.
\newcommand{\gwtitle}[2]{%
  \begin{frame}[plain]
    \begin{tikzpicture}[remember picture,overlay]
      \fill[accent] (current page.south west) rectangle (current page.north east);
      \foreach \i in {1,...,8}{
        \fill[white,opacity=0.08] ([xshift=-\i cm,yshift=\i cm]current page.south east)
          rectangle ++(0.7cm,0.7cm);
      }
    \end{tikzpicture}
    \vfill
    \begin{center}
    {\color{white}\Huge\bfseries #1\par}
    \vspace{6pt}
    {\color{accentlight}\large #2\par}
    \vspace{18pt}
    {\color{accentlight}\small Great Wall \;\textbullet\; Session \gwsessionno{} of 5 \;\textbullet\; hands-on, on Discord\par}
    \end{center}
    \vfill
  \end{frame}}


\begin{document}

\gwtitle{Secrets You Cannot Tell}{Tacit knowledge --- the one secret a wrench can't extract}

% ── Recap + hook ──────────────────────────────────────────────────────────
\begin{frame}{Where we left off}
  \begin{columns}[T]
    \begin{column}{0.6\textwidth}
      Last time: the \textbf{rhetoric trilemma}. Under coercion, comply / deny /
      decoy all lose, because the attacker judges your \emph{words}.
      \vspace{6pt}
      \begin{block}{So the escape must be structural}
        We need a secret you \emph{cannot say}, even if you wanted to --- and a
        constraint the attacker can \emph{verify} is real.
      \end{block}
      Today: the first half --- a secret that isn't words.
    \end{column}
    \begin{column}{0.38\textwidth}
      \begin{tryit}[Check-in]
        \footnotesize One word in \texttt{\#feedback}: what stuck from Session 1?
      \end{tryit}
    \end{column}
  \end{columns}
\end{frame}

% ── Prime intuition ───────────────────────────────────────────────────────
\begin{frame}{A quick thought experiment}
  \begin{columns}[T]
    \begin{column}{0.58\textwidth}
      Describe your \textbf{best friend's face} --- in words --- precisely enough
      that a stranger could pick them out of a crowd of a thousand.
      \vspace{8pt}

      You can't. Yet you'd recognize that face \emph{instantly}.
      \vspace{8pt}

      That gap --- between what you can \emph{do} and what you can \emph{say} ---
      is the entire security model of Great Wall.
    \end{column}
    \begin{column}{0.40\textwidth}
      \memeslot{tacit-face.png}{``describe your friend's face so I can find
      them'' --- brain: \emph{no}.}
    \end{column}
  \end{columns}
\end{frame}

% ── Explicit vs tacit ─────────────────────────────────────────────────────
\begin{frame}{Explicit vs.\ tacit knowledge}
  \begin{columns}[T]
    \begin{column}{0.48\textwidth}
      \begin{block}{Explicit}
        \footnotesize Anything you can \emph{state}: a number, a word, a seed
        phrase. \\[4pt]\textbf{Transmissible} --- which is exactly why it's
        \textbf{coercible}.
      \end{block}
    \end{column}
    \begin{column}{0.48\textwidth}
      \begin{block}{Tacit}
        \footnotesize Knowing-\emph{how} you can't put into words: a face, riding
        a bike, the ``feel'' of a route. \\[4pt]\textbf{Non-transmissible} ---
        so \textbf{non-coercible}.
      \end{block}
    \end{column}
  \end{columns}
  \vspace{8pt}
  \begin{keyidea}
    Recognition $\neq$ recall $\neq$ verbalizable. You can pick your fractal
    point out of the landscape without being able to \emph{read out} its
    coordinates. Coercion can only extract what can be verbalized.
  \end{keyidea}
\end{frame}

\punch{``We can know more\\[6pt] than we can tell.''\\[10pt] {\normalsize\itshape --- Michael Polanyi, 1966}}

% ── Why it defeats the wrench ─────────────────────────────────────────────
\begin{frame}{Why this defeats the wrench}
  A wrench attack raises the cost of \emph{refusing} until you comply. It assumes
  the secret is something you \emph{could} hand over.
  \vspace{6pt}
  \begin{itemize}
    \item There is no sentence you can be tortured into saying that conveys
      ``the point is \emph{there}.''
    \item The only extraction is to sit you at the fractal and have you point ---
      and (Session 4) the fractal you'd point at \textbf{doesn't exist yet} until
      a long computation finishes.
  \end{itemize}
  \vspace{4pt}
  \begin{pitfall}
    Tacit $\neq$ secret-in-the-usual-sense. The attacker may know \emph{everything}
    about the method and that you have a stash (non-obscurity!). What they can't
    get is the non-transmissible \emph{recognition} in your head.
  \end{pitfall}
\end{frame}

% ── TKBA ──────────────────────────────────────────────────────────────────
\begin{frame}{TKBA --- the theory of the whole system}
  Great Wall is \textbf{Tacit Knowledge-Based Authentication}. The four
  properties are \emph{logically coupled}:
  \vspace{4pt}
  \begin{center}\small
  in your head \textbf{(1)} $+$ nobody else's \textbf{(2)} $+$ attacker knows
  \textbf{(3)} \;$\Rightarrow$\; resists coercion \textbf{(4)} \emph{only if} the
  secret is non-transmissible \;$=$\; \textbf{tacit}.
  \end{center}
  \vspace{6pt}
  \begin{keyidea}[The second ingredient]
    TKBA also needs the \emph{interface} to be gated by an \textbf{inescapably
    slow computation} --- otherwise the attacker just says ``try again'' until
    recognition leaks. That gate is Argon2 (Session 4).
  \end{keyidea}
\end{frame}

% ── Design consequence ────────────────────────────────────────────────────
\begin{frame}{A design choice can \emph{break} TKBA}
  \begin{deeper}[Why the area-visualization is hidden by default]
    Anything that turns tacit recall into an \emph{explicit, memorizable fact}
    weakens property 4. So the UX gates the bisection-area rectangles behind a
    debug flag: showing them in normal use would teach you verbalizable facts
    about your own encoding --- facts an attacker could then coerce.
  \end{deeper}
  \vspace{4pt}
  \begin{center}\small
  Security here is a property of \textbf{what you're allowed to know about your
  own secret.}
  \end{center}
  \gwfeedback{Did ``tacit'' actually click, or stay abstract? The pairing
  exercise is the moment of truth.}
\end{frame}

% ── Interaction ───────────────────────────────────────────────────────────
\begin{frame}{Interaction: describe your point}
  \begin{tryit}[Pairs, on Discord --- no screen-share]
    \begin{enumerate}\small
      \item Each of you decodes one guided point in the viewer.
      \item Now \textbf{describe your point's location in words} well enough that
        your partner can click it.
      \item Watch how fast ``up and a bit left, near that flamey bit'' fails to
        pin a 32-bit location.
    \end{enumerate}
  \end{tryit}
  \vspace{4pt}
  The \emph{failure} is the lesson: you just experienced the difference between
  knowing where it is and being able to \emph{say} where it is.
\end{frame}

% ── Homework bridge ───────────────────────────────────────────────────────
\begin{frame}{Before next time}
  \begin{tryit}[Homework into Session 3]
    Keep the viewer handy and read textbook Ch.~3. Next session we stop
    hand-waving about ``encode a secret as a point'' and \textbf{do it} --- and
    watch your exact words round-trip back out.
  \end{tryit}
  \begin{center}\small\color{graytext}
  Question to sit with: if you can't \emph{say} where your point is, how can a
  computer store it \emph{exactly} and get it back? That's Session 3.
  \end{center}
\end{frame}

% ── Check-out ─────────────────────────────────────────────────────────────
\begin{frame}{Check-out --- 60 seconds}
  \begin{columns}[T]
    \begin{column}{0.62\textwidth}
      \begin{block}{In \texttt{\#feedback}}
        \begin{enumerate}\small
          \item Muddiest point today, one line.
          \item Per-session form (pinned) --- especially: \emph{did ``tacit''
            land, or stay abstract?}
        \end{enumerate}
      \end{block}
    \end{column}
    \begin{column}{0.36\textwidth}
      \begin{center}\gwflagicon\\[3pt]
      {\footnotesize\color{attackcol}\itshape Your confusion is the curriculum.}
      \end{center}
    \end{column}
  \end{columns}
\end{frame}

\punch{Next: how a secret becomes\\[6pt] a point on a labyrinth.}

\end{document}
